% ============================================================
\chapter{Conclusiones y trabajo futuro}
\label{cap:conclusiones}
% ============================================================

Este capítulo cierra el trabajo recogiendo las conclusiones vinculadas a cada
objetivo, las consideraciones éticas derivadas del uso de la herramienta, las
limitaciones del estudio y las líneas de trabajo futuro que se desprenden de
los resultados obtenidos. Las conclusiones se presentan en el mismo orden que
los objetivos específicos del Capítulo~\ref{cap:objetivos}, de modo que cada una
quede vinculada explícitamente al objetivo que responde. Las limitaciones y
líneas de trabajo futuro se numeran de forma independiente para facilitar su
referencia cruzada a lo largo del capítulo.

\section{Conclusiones}

Este trabajo ha mostrado la viabilidad de un enfoque híbrido de tres capas para la
detección automática de cadenas de ataque en infraestructuras \textit{Kubernetes}. Las
conclusiones C1, C2, C4, C5, C6 y C7 responden, respectivamente, a los objetivos OE2,
OE4, OE1, OE5, OE3 y OE6 definidos en el Capítulo~\ref{cap:objetivos}; C3 no cierra un
objetivo propio, sino que aporta evidencia complementaria sobre el impacto operacional
del resultado de C2. En conjunto, estas conclusiones cierran el objetivo general
planteado en el Capítulo~\ref{cap:objetivos}. Las conclusiones principales son:

\textbf{C1 — La Capa~1 supera ampliamente el objetivo de diseño (OE2).} El clasificador
\textit{Random Forest} alcanza F1 macro\,=\,0,9946 y \ac{AUC}\,=\,0,9986 en test, superando
el umbral objetivo del 0,85 en 14,5 puntos porcentuales. El modelo de severidad de
tres clases logra F1\,macro\,=\,0,9838, con la clase \textit{high-critical} (96~muestras
de test) en 0,9789 — el resultado más débil de los tres, aunque por encima de 0,95—;
sus 3~errores se clasifican todos como \textit{low-medium}, sin ningún manifiesto
crítico degradado a \textit{clean}.

\textbf{C2 — La Capa~2, en su configuración final, alcanza el objetivo de
Precisión@5 (OE4).} La \ac{GAT} de 3~capas con \textit{embedding} de tipos de
arista, entrenada con entropía cruzada ponderada y selección de
\textit{checkpoints} por Precisión@5 bajo el particionado consciente de familias
sobre el corpus completo de Rahman et~al.\ (599~grafos originales, 37~cadenas),
alcanza P@5\,=\,0,720\,\(\pm\)\,0,160 y F1~macro\,=\,0,803\,\(\pm\)\,0,137
en los 5~pliegues de validación cruzada, superando el objetivo de diseño
P@5\,\(>\)\,0,70 con la semilla desplegada. La lectura honesta de esta cifra
exige dos matices que este trabajo documenta explícitamente: a través de tres
semillas de entrenamiento la media es de 0,68\,\(\pm\)\,0,03 (objetivo dentro
de una desviación típica, con la variación concentrada íntegramente en el
pliegue de la familia \texttt{badpods\_*}), y alcanzar el objetivo exigió
diagnosticar mediante una ablación factorial 2\(\times\)2 que una primera
configuración con \textit{focal loss} y selección por F1 macro
(P@5\,=\,0,480\,\(\pm\)\,0,271) atribuía erróneamente al corpus una regresión
causada por esas dos elecciones (Sección~\ref{sec:ablacion_perdida}). La
ablación de arquitectura, repetida sobre el corpus completo con la configuración
final (Tabla~\ref{tab:arch_ablation}), muestra que la atención sobre tipos de
arista aporta +9,2~puntos de F1 y +4~puntos de P@5 frente a una \ac{GCN} sin
tipos de arista, y que la profundidad (2 frente a 3~capas) no es el factor
determinante del ranking — una propiedad de robustez de la arquitectura.

\textbf{C3 — El \textit{ensemble} final, combinado con la restricción de
viabilidad estructural, mantiene un desempeño operacional sólido sobre el
conjunto de test.} Sobre los 86~grafos de test (5~cadenas de repositorios reales
únicos, sin hermanos de plantilla en entrenamiento; techo P@5\,=\,1,00), el
sistema completo de tres capas alcanza P@1\,=\,1,00, P@3\,=\,1,00 y
P@5\,=\,0,80, recuperando 4~de las 5~cadenas en las primeras 5~posiciones del
ranking sin ningún falso positivo de clúster limpio
(\ac{FPR}\_clean\,=\,0,00); P@1 y P@3 se mantienen idénticos (1,00) a través de
las tres semillas de entrenamiento evaluadas, con P@5 entre 0,80 y 1,00. Parte esencial del resultado es la
restricción de viabilidad estructural (un clúster de un solo manifiesto no puede
albergar una cadena multi-salto, Sección~\ref{sec:gate_estructural}): sin ella
las métricas de ranking caen a P@1\,=\,0,00 y P@5\,=\,0,60, porque 74 de los
86~grafos de test son de un solo nodo. El tamaño reducido del conjunto de test
(intervalo de confianza al 95\,\% de P@5: \([0{,}60,\ 1{,}00]\)) limita la
confianza que puede depositarse en cualquier comparación puntual.

\textbf{C4 — La construcción del corpus y su protocolo de particionado
anti-fuga son la contribución metodológica central (OE1).} El conjunto de datos
final satisface los criterios cuantitativos de OE1: 2.827~manifiestos
(\(\geq\)2.000) y una prevalencia de \texttt{TRUE\_HOST\_PID} del 3,78\,\%
(\(\geq\)2\,\%). La evolución en seis fases — línea base (P@5\,=\,0,40,
82~grafos), aumentación (P@5\,=\,0,52, 277~grafos), corrección
\texttt{HOSTPATH\_MOUNT} (P@5\,=\,0,60, 307~grafos), dataset extendido con
25~cadenas (P@5\,=\,0,88, 471~grafos, protocolo preliminar), corrección
\textit{group-aware} del particionado (P@5\,=\,0,72, mismos 471~grafos) y
corpus completo con familias de plantilla atómicas (P@5\,=\,0,72, 599~grafos,
resultado final) — enseña dos lecciones metodológicas. Primera: la aumentación
de datos exige particionado consciente de grupos, pues parte de la mejora
aparente de la cuarta fase procedía de fuga entre variantes aumentadas y sus
grafos base de evaluación. Segunda, y descubierta en la fase final: los casi
duplicados de plantilla (ocho \texttt{badpods\_*}, nueve \texttt{datadog\_*})
constituyen una segunda vía de fuga cuando se reparten entre particiones; su
corrección hace que las cifras finales midan generalización entre familias no
vistas, una vara de medir más exigente y más honesta que la de las fases~1--4,
cuyas métricas estaban infladas por esta contaminación.

\textbf{C5 — El ensemble cumple los criterios de OE5, con una salvedad
documentada sobre el valor añadido de los pesos optimizados.} La ablación sobre
el conjunto de test (Tabla~\ref{tab:ensemble_ablation}) confirma que el
componente \ac{GNN} es indispensable: la configuración RF-sólo colapsa el
ranking (P@1 de 1,00 a 0,00; P@5 de 0,80 a 0,40) porque el \textit{risk score}
medio por manifiesto de los repositorios reales de cadena es \(\approx\)0,
mientras que \ac{FPR}\_clean\,=\,0,00 se mantiene en todas las configuraciones.
Los criterios de OE5 se cumplen: \ac{FPR}\_clean\,=\,0 y P@5 del ensemble igual
a la del mejor componente individual (0,80). Debe señalarse que los pesos
optimizados por el \ac{GA} no superan a los pesos uniformes ni al \ac{GNN} en
solitario: la función objetivo, escalonada, deja los pesos indeterminados sobre
mesetas amplias, y el conjunto de ajuste \textit{out-of-fold} (dominado por
cadenas de \textit{fixtures} sintéticas) presenta un desplazamiento de
distribución respecto a las cadenas reales de test
(Sección~\ref{sec:layer3}). La aportación del \textit{ensemble} en estos datos
es igualar al mejor componente sin incurrir en los puntos débiles de ninguno,
no mejorar su ranking.

\textbf{C6 — El esquema de cinco tipos de arista amplía la cobertura de escalada de
privilegios (OE3).} El diseño del grafo de clúster incorpora el tipo de arista
\ac{RBAC} mediante detección de roles privilegiados en dos pasadas sobre los
manifiestos, limitando las aristas a cuentas de servicio vinculadas realmente a
roles con permisos de escalada. Esta quinta arista está presente tanto en los
grafos de entrenamiento como en el componente de producción de \textit{kubescan};
sin embargo, la ablación de arquitectura (Tabla~\ref{tab:arch_ablation}) compara
\ac{GAT} frente a \ac{GCN} y no aísla la contribución específica del tipo de arista
\ac{RBAC} frente a los otros cuatro, por lo que su aporte incremental al ranking de
cadenas de ataque queda como trabajo futuro (TF6).

\textbf{C7 — kubescan es una herramienta distribuible, instalable y operacionalmente
viable (OE6).} La evaluación funcional de la Sección~\ref{sec:verificacion_requisitos}
confirma que el sistema se instala sin compiladores externos (\textit{pip install -e
kubescan/}), analiza clústeres de hasta 312~manifiestos en menos de 8~segundos en
hardware de gama media, produce salida en dos formatos (texto y \ac{JSON}) adecuados para
integración en pipelines \acs{CICD}, e incorpora además un modo opcional de consulta
directa a un clúster activo (\texttt{kubescan live}) para auditorías en caliente, sin
que ello altere la vía principal de análisis sobre manifiestos estáticos. La inferencia
es reproducible sobre un mismo \textit{backend} de cómputo mediante la resolución
determinista de \textit{checkpoints} (sujeto a la variación entre \textit{backends}
documentada en L6). El prototipo implementado
en este
trabajo demuestra que el paradigma \textit{shift-left} —detectar cadenas de ataque
antes del despliegue, sin requerir acceso al clúster activo— es técnicamente factible
y operacionalmente eficiente. Complementariamente, la evaluación de usabilidad y
aplicabilidad con usuarios expertos (Sección~\ref{sec:usabilidad}) respalda este
objetivo desde la perspectiva del usuario: la herramienta obtiene una puntuación
\textbf{SUS de 88,3} (valoración «excelente») y una aplicabilidad alta —integración
en CI/CD, priorización de riesgo y confianza en el veredicto (medias de 4,3 sobre
5)—, con una tasa de éxito del 100\,\% en las tareas nucleares de instalación,
escaneo e interpretación (90\,\% global sobre las cinco tareas, incluidas la
extracción \ac{JSON} y la tarea opcional). El principal punto de mejora señalado por
los participantes es la
\textbf{accionabilidad} del informe (A3\,=\,3,7): la herramienta comunica el
riesgo pero no las acciones de remediación ni las reglas concretas tras el
veredicto, lo que se recoge como trabajo futuro (TF8). Estos resultados deben
interpretarse como indicativos por el tamaño de la muestra (\(n=3\),
Sección~\ref{sec:usab_amenazas}).

\section{Consideraciones éticas y uso responsable}
\label{sec:etica}

\textit{kubescan} es una herramienta de uso dual: aunque su propósito es defensivo
—detectar cadenas de ataque potenciales antes de que sean explotadas— la información
que produce (identificación de pods con flags de escape, rutas de privilegio \ac{RBAC})
podría ser utilizada con fines ofensivos si cayera en manos de un actor no autorizado.
En consecuencia, este trabajo adopta los principios de divulgación responsable
(\textit{responsible disclosure}) que rigen la investigación en ciberseguridad.
Los patrones de ataque detectados por \textit{kubescan} se basan exclusivamente en
documentación pública —la guía \ac{NSA}/\ac{CISA} \citep{NSA2022}, el trabajo de BishopFox
\citep{BishopFox2021} y el marco MITRE ATT\&CK for Containers \citep{MITRE2021Containers}—
y no revelan ninguna vulnerabilidad nueva. El código fuente del sistema es de acceso
abierto para facilitar la auditoría independiente de las reglas de detección, lo que
permite a la comunidad verificar que las señales de riesgo generadas son correctas y
no contienen sesgos perjudiciales. El uso de la herramienta sobre clústeres de terceros
sin autorización explícita estaría fuera del ámbito previsto y podría constituir un
acceso no autorizado a sistemas informáticos según la legislación aplicable.

\section{Limitaciones}
\label{sec:limitaciones}

\textbf{L1 — Concentración por familias y varianza entre pliegues.} El
particionado consciente de familias, al mantener cada familia de plantilla como
unidad atómica, concentra las cadenas de validación de cada pliegue en 1--3
familias independientes: el pliegue~0 evalúa íntegramente la familia
\texttt{datadog\_*} (F1\,=\,1,0; P@5\,=\,1,0) y el pliegue~1 la familia
\texttt{badpods\_*} (F1\,=\,0,58), de modo que la varianza entre pliegues
(\(\sigma\)\,=\,0,137 para F1; 0,160 para P@5) mide en gran parte la dispersión
de dificultad entre familias, no inestabilidad del modelo — el análisis
multi-semilla (Sección~\ref{sec:multisemilla}) muestra que cuatro de los cinco
pliegues puntúan de forma idéntica con las tres semillas. El tamaño muestral
efectivo por pliegue (número de familias, no de cadenas) es la limitación
subyacente, y solo puede crecer incorporando repositorios de cadena
estructuralmente independientes (TF1).

\textbf{L2 — Imputación de características extendidas.} Las 7 características
extendidas (basadas en \textit{Checkov}) solo pueden calcularse directamente para
el 47,4\,\% de las filas con \acs{YAML} descargado y se imputan por la mediana de
columna para el resto, diluyendo su señal. La descarga completa de los manifiestos
referenciados en el dataset de Rahman et~al.\ permitiría mejorar estos vectores.
Un experimento de ablación eliminando las 7 características imputadas y
reentrenando únicamente con las 21 restantes de alta cobertura podría cuantificar
el impacto de esta limitación sobre la F1 macro de la Capa~1.

\textbf{L3 — Sesgo de muestreo en el dataset tabular.} Los 10 repositorios más
representados suponen el 47,6\,\% del dataset, y \textit{gitcontroller} por sí solo
aporta el 9,9\,\%. Un esquema de validación cruzada por repositorio ofrecería
estimaciones más conservadoras de la generalización del modelo.

\textbf{L4 — Evaluación en clústeres de producción reales.} Toda la evaluación
se ha realizado sobre repositorios públicos. La validación en entornos de producción
de organizaciones reales es la siguiente etapa crítica para confirmar la validez
externa del sistema.

\textbf{L5 — Etiquetado derivado de reglas sobre el mismo espacio de
características.} Las etiquetas de grafo (clean / isolated / attack\_chain) se
generan mediante una regla determinista definida sobre las mismas \textit{flags} y
aristas que observan los modelos (Sección~\ref{sec:construccion_grafo}). El sistema aprende, por
tanto, una aproximación robusta de dicha regla bajo ruido, enmascaramiento y
observabilidad parcial — no una noción de cadena de ataque independiente de la
regla. Esta circularidad es común en la literatura de detección de
misconfiguraciones, pero limita el alcance de la afirmación «predice cadenas de
ataque». Las vías para romperla son la validación externa frente a herramientas de
grafo de ataque sobre clústeres activos (KubeHound, IceKube), el etiquetado experto
de una submuestra, y la confirmación mediante explotación controlada en entornos de
laboratorio como \textit{kubernetes-goat}.

\textbf{L6 — No determinismo de CUDA en las operaciones de \textit{scatter} del
\ac{GNN} (hallazgo histórico, no reverificado sobre el corpus completo).} Durante
la fase preliminar se observó que, al replicar el entrenamiento de la Capa~2
sobre una GPU NVIDIA A100 (CUDA) en lugar del \textit{backend} MPS/CPU empleado
entonces, las métricas de validación cruzada variaban de forma no trivial (F1
macro entre 0,845 y 0,882 frente a 0,917 en el \textit{backend} de referencia de
esa fase; P@5 entre 0,77 y 0,92 frente a 0,880), manteniendo semilla, datos,
particiones e hiperparámetros idénticos. Estas cifras corresponden a la fase
preliminar de 471~grafos (Tabla~\ref{tab:gnn_evolution}, fase~4, cuyo
\textit{backend} de referencia era MPS/CPU); los resultados finales de esta
memoria se entrenaron y evaluaron sobre GPU NVIDIA T4 (CUDA), por lo que la
amenaza se invierte para el lector que reproduzca en CPU/MPS: es esa réplica la
que puede divergir de las cifras reportadas. El experimento no se ha repetido
sobre el corpus completo, por lo que se reporta como hallazgo histórico sobre el
mecanismo causal, no como una medición vigente sobre el modelo desplegado. Se descartaron como causa la precisión \textit{TF32}
(confirmado en tiempo de ejecución que
\texttt{torch.backends.cuda.matmul.allow\_tf32} permanece desactivada) y la
paralelización del cargador de datos (el conjunto es estático en memoria y el orden
de muestreo está fijado por un generador con semilla). La causa confirmada son las
operaciones \texttt{scatter\_add}/\texttt{scatter\_reduce} empleadas por
\texttt{GATConv} y por las capas de \textit{pooling} global
(\texttt{global\_max\_pool}, \texttt{global\_mean\_pool}), que en CUDA se
implementan mediante \texttt{atomicAdd} sin versión determinista disponible —
confirmado por los propios mantenedores de \textit{PyTorch Geometric} en los
\textit{issues} públicos \#2788 y \#3175 del repositorio. El orden de acumulación
en punto flotante depende de la planificación de hilos de la GPU y, al no ser
la suma en punto flotante asociativa, pequeñas diferencias numéricas cambian el
\textit{ranking} sobre un conjunto de validación reducido (471~grafos, entre 4 y 5
cadenas de ataque por pliegue en ese momento), amplificando el efecto sobre P@5 y
F1 macro. Se incorporaron banderas de determinismo específicas de CUDA
(\texttt{torch.use\_deterministic\_algorithms}, \texttt{cudnn.deterministic}) que,
aunque no resuelven la ausencia de una implementación determinista de
\texttt{scatter\_add}, sí lograron reproducibilidad exacta entre ejecuciones
repetidas sobre la misma configuración de GPU. Esto indica que la brecha
observada no es ruido aleatorio entre ejecuciones sino una divergencia numérica fija
y reproducible entre \textit{backends} (CUDA frente a MPS/CPU); su cierre completo
requeriría sustituir la agregación por la ruta \texttt{SparseTensor} de
\textit{PyTorch Geometric}, documentada como determinista pero no evaluada en este
trabajo, ni reverificada sobre el corpus completo.

\textbf{L7 — Las métricas de las fases previas a la corrección por familias no
son comparables con las finales, y las estimaciones finales llevan intervalos
amplios.} Las evaluaciones de las fases 1--5 estaban infladas por dos vías de
fuga de información corregidas sucesivamente (variantes aumentadas en la fase~5
y familias de plantilla en la fase~6): el modelo «generalizaba» a casi duplicados
de grafos vistos en entrenamiento. Tras la corrección, las cifras finales miden
generalización a familias y repositorios no vistos — una cantidad distinta y
sistemáticamente menor, por lo que ninguna comparación directa entre fases debe
leerse como progresión del modelo. Como segunda cara de la misma limitación, el
conjunto de test contiene solo 5~cadenas (todos los intervalos de confianza de
test tienen una anchura \(\geq\)0,4) y el \textit{pool out-of-fold} contiene
32~cadenas de las que 15 son \textit{fixtures} sintéticas: la evidencia sobre
generalización a cadenas reales descansa en pocas muestras independientes, y
ampliarlas es la vía prioritaria de trabajo futuro (TF1).

\subsection{Amenazas a la validez}
\label{sec:amenazas_validez}

Siguiendo la taxonomía de validez de constructo, interna, de conclusión y
externa empleada en los \textit{ACM SIGSOFT Empirical Standards}, las
limitaciones L1--L7 se organizan del siguiente modo. Esta reorganización no
introduce limitaciones nuevas: reagrupa las ya descritas según el tipo de
amenaza metodológica que representan, para facilitar su lectura conjunta. Cada
categoría remite a las limitaciones L1--L7 correspondientes.

\textbf{Validez de constructo} (¿miden las métricas lo que se pretende medir?):
L6 cuestiona si las métricas reportadas en esta memoria (\textit{backend} CUDA/T4)
coinciden con las que obtendría una réplica en un \textit{backend} CPU/MPS, dado el
no determinismo documentado en las operaciones de \textit{scatter}. L5 cuestiona si la etiqueta de grafo (derivada de reglas sobre
el mismo espacio de características que observan los modelos) constituye un
constructo de «cadena de ataque» independiente de la regla de etiquetado, o una
aproximación circular a ella.

\textbf{Validez interna} (¿las conclusiones causales están justificadas por el
diseño experimental?): L7 impide atribuir diferencias entre fases experimentales
al modelo, porque el protocolo de evaluación cambió entre ellas; la ablación
factorial de la Sección~\ref{sec:ablacion_perdida} es el único contraste interno
con protocolo constante. L2 acota la validez interna de las características
extendidas imputadas mediante \textit{Checkov} frente a las obtenidas
directamente del manifiesto original.

\textbf{Validez de conclusión} (¿son estadísticamente fiables las estimaciones
puntuales reportadas?): el tamaño del conjunto de test (86~grafos, 5~cadenas) y
del \textit{pool out-of-fold} (513~grafos, 32~cadenas) se traduce en los
intervalos de confianza amplios reportados en la Sección~\ref{sec:layer3}
(P@5\,\(\in\,[0{,}60,\ 1{,}00]\) al 95\,\%), que limitan la confianza que puede
depositarse en cualquier comparación puntual entre configuraciones del
\textit{ensemble} (C5). A ello se suma la concentración por familias (L1): el
tamaño muestral efectivo de la validación cruzada es el número de familias
independientes por pliegue (1--3), no el número nominal de cadenas.

\textbf{Validez externa / generalización}: L3 (sesgo de muestreo por
repositorio) y L4 (ausencia de evaluación en clústeres de producción reales)
acotan directamente hasta qué punto los resultados generalizan más allá de los
repositorios públicos utilizados. El desplazamiento de distribución entre el
\textit{pool} de ajuste del \ac{GA} (dominado por \textit{fixtures} sintéticas)
y las cadenas reales de test (Sección~\ref{sec:layer3}) es una amenaza externa
adicional identificada empíricamente: los pesos del \textit{ensemble} ajustados
sobre cadenas sintéticas no aportan ventaja sobre cadenas reales.

\section{Líneas de trabajo futuro}

Los resultados de este trabajo permiten formular las siguientes recomendaciones de
trabajo futuro, ordenadas por impacto esperado. La mayoría de ellas responde
directamente a una limitación de las descritas en la sección anterior, por lo
que cada recomendación indica explícitamente la limitación que atiende cuando
corresponde. Las dos primeras (TF1 y TF2) son las que este trabajo identifica
como de mayor impacto potencial sobre la validez de las métricas reportadas.

\textbf{TF1 — Ampliar la población de cadenas reales estructuralmente
independientes.} Las limitaciones L1 y L7 comparten la misma causa raíz: pocas
familias de cadena independientes. Cada nueva familia de plantilla añade poco
(sus miembros son casi duplicados), mientras que cada repositorio real de cadena
añade una muestra independiente que aumenta el tamaño muestral efectivo de la
validación cruzada, estrecha los intervalos de test y reduce el desplazamiento de
distribución del \textit{pool} de ajuste del \ac{GA}. La incorporación dirigida
de repositorios de infraestructura de organizaciones de código abierto y de
entornos de laboratorio multiclúster —priorizando diversidad estructural sobre
volumen— es la línea de mayor impacto esperado de este trabajo.

\textbf{TF2 — Implementar P@5 global como métrica alternativa.} En lugar de P@5 por
pliegue, evaluar el ranking del modelo sobre el conjunto completo de grafos en un único
ranking global aumentaría el número de candidatos ordenados por evaluación y reduciría
la granularidad escalonada (en pasos de 0,2) de la P@5 por pliegue, que amplifica la
varianza entre pliegues señalada en L1.

\textbf{TF3 — Incorporar análisis semántico de \ac{RBAC} completo.} El parser de \acs{YAML}
actual identifica roles privilegiados por nombre y por verbos de escalada. Una extensión
natural es el análisis de la cadena completa Role/ClusterRole \(\to\) RoleBinding \(\to\)
ServiceAccount \(\to\) Pod para detectar rutas de escalada de privilegios mediante
\ac{RBAC} que no requieran flags de escape directas.

\textbf{TF4 — Validación cruzada dejando un repositorio fuera.} La estrategia
\textit{leave-one-repo-out} ofrece estimaciones más conservadoras de la generalización
que la validación cruzada por fila, y es especialmente relevante dada la concentración
del dataset en pocos repositorios.

\textbf{TF5 — Modo de análisis continuo (\textit{admission controller}).} La integración
de \textit{kubescan} como un \textit{validating admission webhook} de \textit{Kubernetes}
permitiría bloquear el despliegue de manifiestos que eleven el riesgo de cadena de
ataque del clúster en tiempo real, sin necesidad de análisis \textit{post-hoc}.

\textbf{TF6 — Ablación de \textit{pooling} e hiperparámetros restantes.} La
ablación de arquitectura (Tabla~\ref{tab:arch_ablation}) ya cuantifica
empíricamente \ac{GAT} frente a \ac{GCN} y 3~capas frente a 2. Queda pendiente
extender ese mismo protocolo controlado (mismos pliegues, mismas semillas) a la
elección de \textit{mean+max pooling} frente a \textit{mean-only}, así como a un
barrido más amplio de hiperparámetros del optimizador (número de cabezas de
atención, \textit{learning rate}, anchura oculta) que en este trabajo se fijaron
a priori sin búsqueda sistemática (Sección~\ref{sec:capa2_diseno}). Este mismo
protocolo debería extenderse a una ablación por tipo de arista que aísle la
contribución incremental de la arista \ac{RBAC} frente a los otros cuatro tipos,
cerrando la cuestión planteada en C6.

\textbf{TF7 — Ajuste del \textit{ensemble} robusto al desplazamiento de
distribución.} Los pesos del \ac{GA} se ajustan sobre un \textit{pool
out-of-fold} en el que 15 de las 32~cadenas son \textit{fixtures} sintéticas con
\textit{risk scores} \ac{RF} altos, mientras que las cadenas reales de test
tienen \textit{risk score} medio \(\approx\)0; tres refinamientos evaluados
(desempate por \textit{mean reciprocal rank}, agregación del \textit{pool} por
familias, ajuste sobre el \textit{ranking} con restricción estructural) no
corrigieron el sesgo (Sección~\ref{sec:layer3}). Las vías naturales son ajustar
sobre un \textit{pool} restringido a cadenas de repositorios reales (viable a
medida que TF1 amplíe esa población) o sustituir P@5 por una función objetivo
continua menos escalonada, que reduzca las mesetas donde los pesos quedan
indeterminados.

\textbf{TF8 — Explicabilidad y recomendaciones de remediación.} La evaluación
con usuarios expertos (Sección~\ref{sec:usabilidad}) identificó la accionabilidad
del informe como el principal punto de mejora: dos participantes echaron en falta
que la herramienta indicara \textit{qué} corregir y expusiera las reglas concretas
tras la puntuación de riesgo. Como trabajo futuro, el informe podría enriquecerse
con recomendaciones de remediación por manifiesto (por ejemplo, la corrección de
configuración asociada a cada \textit{flag} de escape detectada) y con una
explicación de la contribución de cada capa al veredicto, reforzando la confianza
del operador sin sacrificar la brevedad del \textit{ranking}.
